% Appendices are lettered with Greek capitals, and the figures and tables
% inside them are numbered after their appendix letter.
% The class switches the chapter heading to "ΠΑΡΑΡΤΗΜΑ <γράμμα>. ΤΙΤΛΟΣ"
% once \appendixstartedtrue is set inside the appendix environment.
\begin{appendix}
\appendixstartedtrue
\renewcommand{\thechapter}{\greekalphnum{chapter}}

% Bookmarks and internal anchors cannot carry the Greek numeral, which is not
% expandable inside a PDF string, so they keep a plain arabic counter. Only
% the printed letter changes; every reference still resolves.
\renewcommand{\theHchapter}{appendix.\arabic{chapter}}

\chapter{Πηγαίος κώδικας}
\label{app:code}

Παρατίθενται τα δύο αποσπάσματα στα οποία συμπυκνώνονται οι μεθοδολογικές
επιλογές των κεφαλαίων~\ref{chap:clustering} και~\ref{chap:forecasting-economics}.
Ο πλήρης κώδικας είναι διαθέσιμος στο αποθετήριο της εργασίας. Η συσταδοποίηση
υλοποιείται μέσω της βιβλιοθήκης \textenglish{scikit-learn}
\parencite{pedregosa2011scikit} και η πρόβλεψη χρονοσειρών μέσω της
βιβλιοθήκης \textenglish{statsmodels} \parencite{seabold2010statsmodels}.

\section{Εμφωλευμένη συσταδοποίηση}

Το πρώτο απόσπασμα υλοποιεί τη δομή δένδρου: μια κατάτμηση ως προς την περιοχή
και τον χρόνο και, εντός κάθε συστάδας, δεύτερη κατάτμηση ως προς τα
χαρακτηριστικά RFM. Οι εγγραφές δεύτερου επιπέδου συγκεντρώνονται για ολόκληρη
τη συστάδα πριν την επανάληψη, ώστε το δεύτερο επίπεδο να περιγράφει την αγορά
ως ενιαία οντότητα και όχι κάθε περιοχή χωριστά.

\begin{lstlisting}[language=Python,caption={Συσταδοποίηση δύο επιπέδων},label={lst:nested}]
result = partition(
    rows,
    first_level_features,
    k_range=k_range,
    level=1,
    random_state=random_state,
)

for cluster in result.clusters:
    # Pool every second-level row belonging to any member of this cluster,
    # so the sub-clustering sees the group as a whole rather than one
    # member at a time.
    pooled: list[dict[str, Any]] = []
    for member in cluster.members:
        pooled.extend(second_level_rows.get(member, []))

    if not pooled:
        continue

    child = partition(
        pooled,
        second_level_features,
        k_range=k_range,
        label_prefix=f"{cluster.label}.",
        level=cluster.level + 1,
        random_state=random_state,
    )
    cluster.children = child.clusters
\end{lstlisting}

Η επιλογή του πλήθους συστάδων γίνεται από τα δεδομένα. Για κάθε υποψήφια τιμή
υπολογίζεται ο δείκτης περιγράμματος και κρατείται η καλύτερη, ενώ μια
κατάτμηση που κατέρρευσε σε λιγότερες από δύο ομάδες δεν βαθμολογείται
καθόλου.

\begin{lstlisting}[language=Python,caption={Επιλογή του $k$ με δείκτη περιγράμματος},label={lst:silhouette}]
scaled = StandardScaler().fit_transform(matrix)

for k in candidates:
    model = KMeans(n_clusters=k, random_state=random_state, n_init=10)
    labels = model.fit_predict(scaled)
    # An assignment that collapsed to fewer than k groups cannot be scored.
    if len(set(labels)) < 2:
        continue
    score = float(silhouette_score(scaled, labels))
    result.scores[k] = score
    if best_score is None or score > best_score:
        best_k, best_score, best_labels = k, score, labels
\end{lstlisting}

\section{Αποσύνθεση των εσόδων σε ροές}

Το δεύτερο απόσπασμα κατατάσσει τα έσοδα κάθε πελάτη σε μία ακριβώς από τις
έξι ροές της ενότητας~\ref{sec:growth-accounting}. Η αποκλειστικότητα αυτή
είναι που καθιστά τις λογιστικές ταυτότητες ελέγξιμες.

\begin{lstlisting}[language=Python,caption={Κατάταξη των εσόδων σε έξι ροές},label={lst:flows}]
for customer, revenue in current.items():
    before = previous.get(customer, 0.0)
    if before > 0:
        # Active in both periods: the overlap is retained, and the
        # difference is expansion or contraction.
        result.retained += min(before, revenue)
        if revenue > before:
            result.expansion += revenue - before
        else:
            result.contraction += before - revenue
    elif customer in seen_before:
        result.resurrected += revenue
    else:
        result.new += revenue

for customer, before in previous.items():
    if current.get(customer, 0.0) == 0 and before > 0:
        result.churned += before
\end{lstlisting}

\chapter{Αναλυτικά αποτελέσματα}
\label{app:results}

\section{Πλήρες πλέγμα ευαισθησίας των μοναδιαίων οικονομικών}

Ο πίνακας~\ref{tab:sensitivity} του κεφαλαίου~\ref{chap:forecasting-economics}
παραθέτει επιλεγμένα κελιά. Εδώ παρατίθεται το πλέγμα στο σύνολό του, ώστε η
κατανομή των θετικών και αρνητικών αποτελεσμάτων να είναι ελέγξιμη και όχι
επιλεγμένη.

Η μέση παρατηρούμενη αξία ανά πελάτη είναι 2.895. Η αξία διάρκειας ζωής
υπολογίζεται επί του μεικτού κέρδους, δηλαδή ως γινόμενο της αξίας αυτής με το
υποτιθέμενο περιθώριο.

\begin{table}[H]
\centering
\caption{Πλήρες πλέγμα ευαισθησίας, 16 κελιά}
\label{tab:sensitivity-full}
\begin{tabular}{rrrrc}
\toprule
Περιθώριο & CAC & LTV & LTV:CAC & Πληροί 3:1 \\
\midrule
0,15 & 100 &   434 & 4,34 & ναι \\
0,15 & 250 &   434 & 1,74 & όχι \\
0,15 & 500 &   434 & 0,87 & όχι \\
0,15 & 900 &   434 & 0,48 & όχι \\
\midrule
0,30 & 100 &   868 & 8,68 & ναι \\
0,30 & 250 &   868 & 3,47 & ναι \\
0,30 & 500 &   868 & 1,74 & όχι \\
0,30 & 900 &   868 & 0,96 & όχι \\
\midrule
0,45 & 100 & 1.303 & 13,03 & ναι \\
0,45 & 250 & 1.303 &  5,21 & ναι \\
0,45 & 500 & 1.303 &  2,61 & όχι \\
0,45 & 900 & 1.303 &  1,45 & όχι \\
\midrule
0,60 & 100 & 1.737 & 17,37 & ναι \\
0,60 & 250 & 1.737 &  6,95 & ναι \\
0,60 & 500 & 1.737 &  3,47 & ναι \\
0,60 & 900 & 1.737 &  1,93 & όχι \\
\bottomrule
\end{tabular}
\end{table}

Οκτώ κελιά πληρούν το σημείο αναφοράς 3:1 και οκτώ όχι. Το κρίσιμο σημείο δεν
είναι η αναλογία αυτή, η οποία εξαρτάται από το επιλεγμένο εύρος, αλλά το ότι
το συμπέρασμα αντιστρέφεται εντός εύρους παραδοχών που όλες θα μπορούσαν να
χαρακτηριστούν εύλογες.

\section{Δείκτες περιγράμματος ανά υποψήφιο πλήθος συστάδων}

Οι βαθμολογίες του πρώτου επιπέδου παρατίθενται στον
πίνακα~\ref{tab:silhouette}. Η τιμή $k=2$ υπερέχει σαφώς, και η πτώση από
0,8793 σε 0,7206 με μία μόνο επιπλέον συστάδα δείχνει ότι η δομή των δεδομένων
είναι δυαδική και όχι βαθμιαία: η εγχώρια αγορά έναντι των εξαγωγικών.

\end{appendix}
